\section{Background}
\label{sec:background}

\paragraph{Agents.} In artificial intelligence and cognitive science, an agent is a system that perceives its environment and acts upon it \citep{russell2020aima}, usually in a repeating cycle of perception, decision and action, supported by memory. Cognitive science distinguishes a short-lived \emph{working memory} from long-term memories that are \emph{episodic} (what happened), \emph{semantic} (what is true) and \emph{procedural} (how to do things).

\paragraph{Language agents.} The loop that harnesses implement goes back to ReAct \citep{yao2023react}, which interleaves a model's reasoning with actions whose results are returned to it as observations. Later systems elaborated the loop with patterns that, in our terms, live in particular primitives. Reflexion \citep{shinn2023reflexion} has the agent write reflections on failed attempts and carries them into the next attempt: a decision in \prim{Control flow}, which runs the attempts, and in \prim{Context}, which carries the reflections forward. MemGPT \citep{packer2023memgpt} manages what is kept inside and outside a limited context window, paging information in and out as needed: a decision in \prim{Context}.

\paragraph{CoALA.} \citet{sumers2024coala} proposed \emph{Cognitive Architectures for Language Agents}, which organizes language agents by their memory (working, episodic, semantic and procedural), their action space (external \emph{grounding} actions and internal \emph{retrieval}, \emph{reasoning} and \emph{learning} actions), and a decision cycle of planning and execution. CoALA describes \emph{what an agent is}, as a cognitive architecture. Our concern is narrower and more operational: \emph{what software an agent is made of}, as a system an engineer builds. The two views are complementary, and \S\ref{sec:primitives} maps between them.

\paragraph{Agent-computer interfaces.} \citet{yang2024sweagent} showed that the interface through which a model acts on a computer (the commands it is given, their output format, and guardrails such as linting edits) substantially changes performance on software engineering tasks, and named this the \emph{agent-computer interface}. In our terms, their interface is a set of decisions within the \prim{Output} and \prim{Context} primitives.

\paragraph{Practitioner guidance.} \citet{anthropic2024effective} distinguishes \emph{workflows}, in which code fixes the sequence of model calls, from \emph{agents}, in which the model directs its own process and tool use, and describes composable patterns such as prompt chaining, routing and orchestrator--worker designs. We treat both workflows and agents as harnesses that differ in their \prim{Control flow} (\S\ref{sec:casestudy}).
